\chapter{电路：回路、串联与并联}\label{ch:g7:circuits-series-parallel}

四年的电路学习已为你备好工具箱：回路、\omterm{def:g6:circuits-and-safety:diagram}{电路图}，
以及两种重要的接法。但有一个角色一直默默无闻地工作，
它就是那个流动的东西。今年，它终于有了名字，
还有它的方向，以及它最奇特的性质：它永远不会被用完。
\omterm{def:g7:circuits-series-parallel:current}{电流}登上舞台后，\omterm{def:g4:series-circuits:series}{串联}与\omterm{def:g5:series-parallel-circuits:parallel}{并联}便不再只是配方，
而成为规律。

\section{电流登场}

\begin{definition}[电流]\label{def:g7:circuits-series-parallel:current}
\emph{电流}\index{电流}是指沿\omterm{def:g3:first-electric-circuit:circuit}{闭合电路}流动的
一群带电粒子：数量多得难以想象，个头小得难以想象，
与\omterm{def:g7:states-of-matter:molecule}{分子}世界颇为相近——
它们在\omterm{def:g3:first-electric-circuit:battery}{电池}的推动下穿过导线中的金属，整齐地前进。
没有回路，便没有前进：电流只在
通路闭合时才存在。（运动的粒子究竟是什么，
金属又如何让它们通过，是一个精彩的故事，
留待本书的最后几年再讲。）
\end{definition}

\begin{proposition}[约定的方向]\label{prop:g7:circuits-series-parallel:direction}
按照国际通行的约定，\omterm{def:g6:circuits-and-safety:diagram}{电路图}中的\omterm{def:g7:circuits-series-parallel:current}{电流}方向标注为
\emph{在\omterm{def:g3:first-electric-circuit:battery}{电池}外部，从 $+$ \omterm{def:g3:first-electric-circuit:battery}{接线柱}出发，绕电路一周，到达 $-$ \omterm{def:g3:first-electric-circuit:battery}{接线柱}}。
世界上每张\omterm{def:g6:circuits-and-safety:diagram}{电路图}上的箭头都遵循这一约定，
这一约定早在人们还无法询问运动粒子本身之前就已确立。
（当物理学终于能够询问时，
答案里藏着一个著名的小笑话，
留待以后某一年再讲。）
\end{proposition}

\begin{omfigure}
\begin{tikzpicture}[scale=0.9]
  \draw (0,0) to[battery1, l=电池] (0,3.0)
    to[short] (2.0,3.0)
    to[lamp] (4.4,3.0)
    to[short] (6.4,3.0)
    to[short] (6.4,0)
    to[short] (0,0);
  \draw[-Stealth, very thick, omThm] (1.1,3.25) -- (1.9,3.25);
  \draw[-Stealth, very thick, omThm] (6.65,1.9) -- (6.65,1.1);
  \draw[-Stealth, very thick, omThm] (5.3,-0.25) -- (4.5,-0.25);
  \draw[-Stealth, very thick, omThm] (1.9,-0.25) -- (1.1,-0.25);
  \node[right, font=\small, omThm] at (7.0,1.5)
    {电流：从 $+$ 流出，绕行，流入 $-$};
\end{tikzpicture}

{\small 全世界电工共同遵守的约定：\omterm{def:g7:circuits-series-parallel:current}{电流}
箭头从\omterm{def:g3:first-electric-circuit:battery}{电池}的 $+$ 极流出，绕行一周，
由 $-$ 极流回。}
\end{omfigure}

\begin{proposition}[电流不会被消耗]\label{prop:g7:circuits-series-parallel:notconsumed}
用电器并不会“吃掉”\omterm{def:g7:circuits-series-parallel:current}{电流}。进入\omterm{def:g3:first-electric-circuit:bulb}{灯泡}、
\omterm{ex:g6:circuits-and-safety:motor}{电动机}或\omterm{ex:g6:circuits-and-safety:motor}{蜂鸣器}的\omterm{def:g7:circuits-series-parallel:current}{电流}，会原封不动地继续流走：
用电器取走的是来自\omterm{def:g3:first-electric-circuit:battery}{电池}储存的\emph{\omterm{def:g5:everyday-energy:energy}{能量}}，由流动的\omterm{def:g7:circuits-series-parallel:current}{电流}输送。
因此，在任何\omterm{def:g4:series-circuits:series}{串联}回路中，同一股\omterm{def:g7:circuits-series-parallel:current}{电流}流过每一个用电器，
回到\omterm{def:g3:first-electric-circuit:battery}{电池}的\omterm{def:g7:circuits-series-parallel:current}{电流}与离开时
一样多。
\end{proposition}

\begin{example}[对错误观念的检验]\label{ex:g7:circuits-series-parallel:trial}
易犯的错误是：“第一个\omterm{def:g3:first-electric-circuit:bulb}{灯泡}耗掉了一些\omterm{def:g7:circuits-series-parallel:current}{电流}，所以
第二个\omterm{def:g3:first-electric-circuit:bulb}{灯泡}必定更暗。”电路本身即可驳斥这一说法。把
两个\emph{相同的}\omterm{def:g3:first-electric-circuit:bulb}{灯泡}\omterm{def:g4:series-circuits:series}{串联}起来，仔细观察：两个\omterm{def:g3:first-electric-circuit:bulb}{灯泡}亮度完全相同，
对调它们的位置，结果毫无改变。
如果\omterm{def:g7:circuits-series-parallel:current}{电流}在途中被消耗，第二个\omterm{def:g3:first-electric-circuit:bulb}{灯泡}就应当总是较暗的那个，
无论哪个\omterm{def:g3:first-electric-circuit:bulb}{灯泡}放在那里。事实并非如此。\omterm{def:g7:circuits-series-parallel:current}{电流}
完整地到达；两个\omterm{def:g3:first-electric-circuit:bulb}{灯泡}共同分享的是\omterm{def:g3:first-electric-circuit:battery}{电池}的\emph{推动力}，
这正是两个\omterm{def:g3:first-electric-circuit:bulb}{灯泡}都比单独一个\omterm{def:g3:first-electric-circuit:bulb}{灯泡}暗的原因，而且暗的程度相同。
\end{example}

\section{两类电路的规律}

\begin{proposition}[串联规律]\label{prop:g7:circuits-series-parallel:series}
在\omterm{def:g4:series-circuits:series}{串联电路}——即单一回路——中，同一股\omterm{def:g7:circuits-series-parallel:current}{电流}
流经每一个用电器：一条通路，一股\omterm{def:g7:circuits-series-parallel:current}{电流}，没有例外。
由此可得以下结论：任何一处断开，整个回路的\omterm{def:g7:circuits-series-parallel:current}{电流}随即中断；
相同的用电器平分工作，无论它们处于何处；
用电器在回路中的排列次序无关紧要，
因为\omterm{def:g7:circuits-series-parallel:current}{电流}同样地流经它们全部。
\end{proposition}

\begin{proposition}[并联规律]\label{prop:g7:circuits-series-parallel:parallel}
在\omterm{def:g5:series-parallel-circuits:parallel}{节点}处，\omterm{def:g7:circuits-series-parallel:current}{电流}\emph{分流}：一部分\omterm{def:g7:circuits-series-parallel:current}{电流}
进入各条\omterm{def:g5:series-parallel-circuits:parallel}{支路}，各部分在远端的\omterm{def:g5:series-parallel-circuits:parallel}{节点}处重新汇合，
既不减少，也不增加。每条\omterm{def:g5:series-parallel-circuits:parallel}{支路}都得到\omterm{def:g3:first-electric-circuit:battery}{电池}的
完整推动力，因此每个用电器都如同单独工作；
某条\omterm{def:g5:series-parallel-circuits:parallel}{支路}被阻断时，它便无\omterm{def:g7:circuits-series-parallel:current}{电流}通过，而其他\omterm{def:g5:series-parallel-circuits:parallel}{支路}
依然照常工作。
\end{proposition}

\begin{omfigure}
\begin{tikzpicture}[scale=0.85]
  \draw (0,0) to[battery1] (0,3.2)
    to[short] (2.4,3.2);
  \draw (2.4,3.2) to[short] (2.4,2.3)
    to[lamp] (5.6,2.3) to[short] (5.6,3.2);
  \draw (2.4,3.2) to[short] (2.4,4.1)
    to[lamp] (5.6,4.1) to[short] (5.6,3.2);
  \draw (5.6,3.2) to[short] (7.2,3.2) to[short] (7.2,0)
    to[short] (0,0);
  \fill (2.4,3.2) circle (2pt);
  \fill (5.6,3.2) circle (2pt);
  \draw[-Stealth, very thick, omThm] (1.0,3.45) -- (1.9,3.45);
  \draw[-Stealth, thick, omThm] (2.85,4.35) -- (3.65,4.35);
  \draw[-Stealth, thick, omThm] (2.85,2.05) -- (3.65,2.05);
  \draw[-Stealth, very thick, omThm] (6.05,3.45) -- (6.95,3.45);
  \node[above, font=\small, omThm] at (1.45,3.55) {整股电流};
  \node[above, font=\small, omThm] at (6.5,3.55) {依然完整};
\end{tikzpicture}

{\small \omterm{def:g5:series-parallel-circuits:parallel}{并联}定律图示：队伍在第一个\omterm{def:g5:series-parallel-circuits:parallel}{节点}处分开，
每一部分沿自己的\omterm{def:g5:series-parallel-circuits:parallel}{支路}前行一圈，
然后在第二个\omterm{def:g5:series-parallel-circuits:parallel}{节点}处重新汇合。}
\end{omfigure}

\begin{example}[借助电工的眼光读亮度]\label{ex:g7:circuits-series-parallel:brightness}
一节\omterm{def:g3:first-electric-circuit:battery}{电池}，相同的\omterm{def:g3:first-electric-circuit:bulb}{灯泡}。单独一只时：满亮——
作为参照。两只\omterm{def:g4:series-circuits:series}{串联}：两者同样暗——一股队伍，共享推力。
两只\omterm{def:g5:series-parallel-circuits:parallel}{并联}：两者都很亮——每条\omterm{def:g5:series-parallel-circuits:parallel}{支路}都享有
整份推力（\omterm{def:g3:first-electric-circuit:battery}{电池}耗尽的\omterm{def:g5:speed-distance-time:speed}{速度}约快一倍：
没有什么是白来的）。三只\omterm{def:g4:series-circuits:series}{串联}：更暗，三者亮度相同。电工一瞥即知：
看回路同伴数，便知暗淡程度；数\omterm{def:g5:series-parallel-circuits:parallel}{支路}数，
便可算出\omterm{def:g3:first-electric-circuit:battery}{电池}的“账单”。
\end{example}

\begin{method}[预判任一灯泡的命运]\label{met:g7:circuits-series-parallel:predict}
给定\omterm{def:g6:circuits-and-safety:diagram}{电路图}，对每个\omterm{def:g3:first-electric-circuit:bulb}{灯泡}：
\begin{enumerate}
  \item \emph{存在性}：是否至少存在一条无断路的
        从 $+$ 经此\omterm{def:g3:first-electric-circuit:bulb}{灯泡}至 $-$？无回路，则无光。
  \item \emph{路上同伴}：哪些器件共享
        此\omterm{def:g3:first-electric-circuit:bulb}{灯泡}经过的\emph{每一条}回路？它们即是其\omterm{def:g4:series-circuits:series}{串联}
        同伴——同伴越多，光越暗。
  \item \emph{\omterm{def:g5:series-parallel-circuits:parallel}{支路}邻居}：哪些器件位于其他
        \omterm{def:g5:series-parallel-circuits:parallel}{支路}上？它们既不使此灯变暗，也不与之同时变暗。
  \item \emph{\omterm{ex:g3:first-electric-circuit:switch}{开关}}：断开的\omterm{ex:g3:first-electric-circuit:switch}{开关}恰好使其所在的
        回路全部失效——将其移去后重新执行步骤 1。
\end{enumerate}
\end{method}

\begin{remark}[这一误解为何根深蒂固]\label{rem:g7:circuits-series-parallel:whyhard}
“必定有东西被用掉了——\omterm{def:g3:first-electric-circuit:battery}{电池}会耗尽！”没错：
\omterm{def:g3:first-electric-circuit:battery}{电池}所储存的\emph{\omterm{def:g5:everyday-energy:energy}{能量}}被耗用，经流动的队伍交给\omterm{def:g3:first-electric-circuit:bulb}{灯泡}，化为光
与热。但队伍中的成员本身
毫发无损地循环，就像缆车吊椅日复一日
把滑雪者送上山坡，而吊椅本身从未被消耗。请将两笔账分开：
\omterm{def:g7:circuits-series-parallel:current}{电流}是循环的，\omterm{def:g5:everyday-energy:energy}{能量}是被输送的——
如此，本书中的每个电路都经得起检验。用来计数队伍本身的仪器
明年才会与我们相见。
\end{remark}

\section{练习}

\begin{exercise}[$\star$]\label{exo:g7:circuits-series-parallel:1}
什么是\omterm{def:g7:circuits-series-parallel:current}{电流}？\omterm{def:g7:circuits-series-parallel:current}{电流}在电路中何时存在，
何时不存在？
\end{exercise}

\begin{exercise}[$\star$]\label{exo:g7:circuits-series-parallel:2}
指出\omterm{def:g7:circuits-series-parallel:current}{电流}的规定方向。在单\omterm{def:g3:first-electric-circuit:bulb}{灯泡}回路的\omterm{def:g6:circuits-and-safety:diagram}{电路图}上，
箭头指向哪里？
\end{exercise}

\begin{exercise}[$\star$]\label{exo:g7:circuits-series-parallel:3}
两个相同的\omterm{def:g3:first-electric-circuit:bulb}{灯泡}\omterm{def:g4:series-circuits:series}{串联}时，亮度
完全相等。这一现象驳斥了哪种流行的错误观点，又如何驳斥？
\end{exercise}

\begin{exercise}[$\star$]\label{exo:g7:circuits-series-parallel:4}
在含\omterm{ex:g6:circuits-and-safety:motor}{电动机}与\omterm{def:g3:first-electric-circuit:bulb}{灯泡}的\omterm{def:g4:series-circuits:series}{串联}回路中，流过\omterm{ex:g6:circuits-and-safety:motor}{电动机}的
\omterm{def:g7:circuits-series-parallel:current}{电流}与流过\omterm{def:g3:first-electric-circuit:bulb}{灯泡}的\omterm{def:g7:circuits-series-parallel:current}{电流}相比较。\omterm{def:g4:series-circuits:series}{串联}定律
说明了什么？
\end{exercise}

\begin{exercise}[$\star$]\label{exo:g7:circuits-series-parallel:5}
\omterm{def:g5:series-parallel-circuits:parallel}{并联电路}的第一个\omterm{def:g5:series-parallel-circuits:parallel}{节点}处，\omterm{def:g7:circuits-series-parallel:current}{电流}会
发生怎样的变化？第二个\omterm{def:g5:series-parallel-circuits:parallel}{节点}处呢？各\omterm{def:g5:series-parallel-circuits:parallel}{支路}所分得的
\omterm{def:g3:first-electric-circuit:battery}{电池}推力有何规律？
\end{exercise}

\begin{exercise}[$\star$]\label{exo:g7:circuits-series-parallel:6}
一节\omterm{def:g3:first-electric-circuit:battery}{电池}，相同的\omterm{def:g3:first-electric-circuit:bulb}{灯泡}：将一个\omterm{def:g3:first-electric-circuit:bulb}{灯泡}的亮度按下列情形排序——
单独时；\ 有一个\omterm{def:g4:series-circuits:series}{串联}同伴时；\ 有一个\omterm{def:g5:series-parallel-circuits:parallel}{并联}
邻居时；\ 有两个\omterm{def:g4:series-circuits:series}{串联}同伴时。
\end{exercise}

\begin{exercise}[$\star$]\label{exo:g7:circuits-series-parallel:7}
哪种接法使\omterm{def:g3:first-electric-circuit:battery}{电池}耗尽得更快：两个\omterm{def:g3:first-electric-circuit:bulb}{灯泡}\omterm{def:g4:series-circuits:series}{串联}，
还是同样两个\omterm{def:g3:first-electric-circuit:bulb}{灯泡}\omterm{def:g5:series-parallel-circuits:parallel}{并联}？请用队伍的语言说明原因。
\end{exercise}

\begin{exercise}[$\star$]\label{exo:g7:circuits-series-parallel:8}
解释
\cref{rem:g7:circuits-series-parallel:whyhard}中的缆车图景：什么充当
吊椅，什么充当滑雪者，真正被用掉的又是什么？
\end{exercise}

\begin{exercise}[$\star\star$]\label{exo:g7:circuits-series-parallel:9}
一张\omterm{def:g6:circuits-and-safety:diagram}{电路图}：\omterm{def:g3:first-electric-circuit:battery}{电池}、\omterm{ex:g3:first-electric-circuit:switch}{开关}，然后是\omterm{def:g3:first-electric-circuit:bulb}{灯泡} A；\omterm{def:g3:first-electric-circuit:bulb}{灯泡} A 之后有一个\omterm{def:g5:series-parallel-circuits:parallel}{节点}
分出两条\omterm{def:g5:series-parallel-circuits:parallel}{支路}——一条接\omterm{def:g3:first-electric-circuit:bulb}{灯泡} B，另一条接\omterm{def:g3:first-electric-circuit:bulb}{灯泡} C——
两条\omterm{def:g5:series-parallel-circuits:parallel}{支路}在回到\omterm{def:g3:first-electric-circuit:battery}{电池}前汇合。利用
\cref{met:g7:circuits-series-parallel:predict}：A 的\omterm{def:g4:series-circuits:series}{串联}同伴
有哪些？C 所在\omterm{def:g5:series-parallel-circuits:parallel}{支路}被切断时，A 的亮度会怎样？
\end{exercise}

\begin{exercise}[$\star\star$]\label{exo:g7:circuits-series-parallel:10}
装饰灯需要一个设计方案：十个\omterm{def:g3:first-electric-circuit:bulb}{灯泡}都以
满亮度、彼此独立地由同一电源供电——但整组灯须
服从一个总\omterm{ex:g3:first-electric-circuit:switch}{开关}。请画出或描述
接线方式，并引用两条定律。
\end{exercise}

\begin{exercise}[$\star\star$]\label{exo:g7:circuits-series-parallel:11}
在两个\omterm{def:g3:first-electric-circuit:bulb}{灯泡}的\omterm{def:g5:series-parallel-circuits:parallel}{并联电路}中，\omterm{def:g3:first-electric-circuit:bulb}{灯泡} 1 所在\omterm{def:g5:series-parallel-circuits:parallel}{支路}
为粗铜线，而\omterm{def:g3:first-electric-circuit:bulb}{灯泡} 2 所在\omterm{def:g5:series-parallel-circuits:parallel}{支路}则包含一段长而细的导线。
\omterm{def:g3:first-electric-circuit:bulb}{灯泡} 2 的亮度略弱。这说明，\omterm{def:g5:series-parallel-circuits:parallel}{支路}本身的
“路”如何能够限制其所分得的队伍份额？（下一
学年将为这种限制命名，并给出其定律。）
\end{exercise}

\begin{exercise}[$\star\star\star$]\label{exo:g7:circuits-series-parallel:12}
一个\omterm{def:g3:first-electric-circuit:battery}{电池}为两条\omterm{def:g5:series-parallel-circuits:parallel}{并联}\omterm{def:g5:series-parallel-circuits:parallel}{支路}供电：\omterm{def:g5:series-parallel-circuits:parallel}{支路}一仅含\omterm{def:g3:first-electric-circuit:bulb}{灯泡} X，
\omterm{def:g5:series-parallel-circuits:parallel}{支路}二则含\omterm{def:g3:first-electric-circuit:bulb}{灯泡} Y 与 Z，二者\omterm{def:g4:series-circuits:series}{串联}。请依次比较三只\omterm{def:g3:first-electric-circuit:bulb}{灯泡}的亮度，
并分别用两条定律说明理由。
随后预测若\omterm{def:g3:first-electric-circuit:bulb}{灯泡} Z 烧坏，每只\omterm{def:g3:first-electric-circuit:bulb}{灯泡}的状况；
若改为\omterm{def:g3:first-electric-circuit:bulb}{灯泡} X 烧坏，又会如何。
\end{exercise}

\section{问题：灯塔看守人的控制面板}

\begin{problem}[{周末问题——重新布线旧灯塔；
一组电池，多项任务；看守人的
故障之夜}]\label{pb:g7:circuits-series-parallel:1}
灯塔的\omterm{def:g3:first-electric-circuit:battery}{电池}组须为以下用电器供电：大灯、
雾笛（一种强力\omterm{ex:g6:circuits-and-safety:motor}{蜂鸣器}）、看守人的阅读灯以及
楼梯灯，并须经受看守人排障的
一整夜考验。

\medskip\noindent\textbf{第一部分——设计。}
\begin{enumerate}
  \item 大灯决不可因其他用电器接通而变暗。
        整套装置须采用何种接线？
        为何任何\omterm{def:g4:series-circuits:series}{串联}都会
        使船只陷入险境？
  \item 每项用电器都需要各自的\omterm{ex:g3:first-electric-circuit:switch}{开关}，看守人又希望
        设一个用于暴风雨时的总\omterm{ex:g3:first-electric-circuit:switch}{开关}。请布置全部五个\omterm{ex:g3:first-electric-circuit:switch}{开关}。
  \item 雾笛与楼梯灯决不可
        因彼此的故障而无法工作。这一点
        是否已经得到保证？依据哪条定律？
  \item 画出或描述完整的\omterm{def:g6:circuits-and-safety:diagram}{电路图}：\omterm{def:g3:first-electric-circuit:battery}{电池}组、总\omterm{ex:g3:first-electric-circuit:switch}{开关}、
        四条\omterm{def:g5:series-parallel-circuits:parallel}{支路}及四个\omterm{def:g5:series-parallel-circuits:parallel}{支路}\omterm{ex:g3:first-electric-circuit:switch}{开关}。
\end{enumerate}

\medskip\noindent\textbf{第二部分——看守人的疑问。}
\begin{enumerate}[resume]
  \item 看守人担心：“四项用电器同时取电，
        每条\omterm{def:g5:series-parallel-circuits:parallel}{支路}是否仍能获得\omterm{def:g3:first-electric-circuit:battery}{电池}组的全部推力？”
        请用\omterm{def:g5:series-parallel-circuits:parallel}{并联}定律作答。
  \item “\omterm{def:g7:circuits-series-parallel:current}{电流}流经大灯时是否会逐渐耗尽？”
        请为看守人解惑，
        如有帮助，可借助缆车作比喻。
  \item 大雾之夜，所有用电器同时工作。
        \omterm{def:g5:series-parallel-circuits:parallel}{并联}的“慷慨”要付出什么代价？
        这一代价在灯塔的哪一部分承担？
  \item 只点亮阅读灯。看守人习惯于\omterm{def:g6:measurement-in-science:measuring}{测量}，
        请追踪\omterm{def:g7:circuits-series-parallel:current}{电流}的整个回路：
        说明其依次经过的每个元件，
        以及它所绕开的每个元件。
\end{enumerate}

\medskip\noindent\textbf{第三部分——故障之夜。}
\begin{enumerate}[resume]
  \item 午夜：楼梯\omterm{def:g3:first-electric-circuit:bulb}{灯泡}烧坏了。还有哪些灯
        会熄灭？依据哪条定律？
  \item 凌晨一点：笨拙的修理留下一段裸露导线，
        直接连接\omterm{def:g3:first-electric-circuit:battery}{电池}组的两个\omterm{def:g3:first-electric-circuit:battery}{接线柱}，位于所有\omterm{def:g5:series-parallel-circuits:parallel}{支路}的
        上游。请命名这一事件，并解释为何灯塔的
        每项用电器此刻都同时失效——
        \omterm{def:g7:circuits-series-parallel:current}{电流}去了哪里？
  \item 凌晨两点，故障已排除：看守人希望为大灯加以保护，
        使此类\omterm{def:g6:circuits-and-safety:short}{短路}不能再次使其熄灭。现有一组独立的
        第二\omterm{def:g3:first-electric-circuit:battery}{电池}组可供使用。
        请提出接线方案。
  \item 黎明：请写下看守人的值班日志——三句话：
        今夜哪条定律保障了船只安全，
        哪一种错误观念被这一夜推翻，
        以及哪一根导线造成了最糟糕的一\omterm{ex:g2:measuring-time:clock}{小时}。
\end{enumerate}
\end{problem}
